\section{Turning Points: Stalingrad and the Global War}

\subsection{Stalingrad 1942-1943: The Encirclement of the Sixth Army}
After the failure to capture Moscow in 1941, Germany launched a renewed offensive in 1942 aimed at the oilfields of the Caucasus and the industrial city of Stalingrad on the Volga. The city carried symbolic weight bearing Stalin's name, and both leaders invested it with prestige beyond its strategic value. German forces drove deep into Soviet territory through the summer, but overextended supply lines and dispersed objectives strained the offensive. By autumn the Sixth Army under General Paulus had become entangled in grinding urban combat that consumed men and material at a terrible rate.

The battle within Stalingrad degenerated into ferocious house-to-house fighting amid ruins, factories, and rubble. Soviet defenders clung to narrow positions along the river, supplied across the Volga under constant fire, trading space and lives to hold the city. The Germans captured most of the urban area but never fully eliminated resistance. This attritional struggle fixed the bulk of German strength in a confined zone while Soviet commanders prepared a far larger operation against the weaker flanks held by Romanian, Hungarian, and Italian allies surrounding the city.

In November 1942 the Red Army launched Operation Uranus, striking the vulnerable flanks and rapidly encircling the Sixth Army within days. Hitler forbade any breakout, insisting the army hold its position and rely on aerial resupply that proved hopelessly inadequate. Trapped, starving, and freezing through the winter, the German troops endured catastrophic conditions while relief efforts failed. The encirclement transformed an offensive into a disaster, demonstrating both the growing operational skill of Soviet commanders and the rigid, ideologically driven decision-making at the summit of the German war effort.

Paulus surrendered the remnants of his army in early February 1943, ending one of the largest and bloodiest battles in human history. The Wehrmacht lost roughly a quarter of a million men killed or captured, along with vast quantities of irreplaceable equipment. Beevor describes the catastrophe as a profound psychological as well as military blow, shattering the myth of German invincibility. The strategic initiative on the Eastern Front passed permanently to the Soviet Union, which would steadily push German forces westward over the following two years.

Stalingrad's significance extended beyond casualty figures to the broader trajectory of the war. It marked the moment when Germany lost the capacity to dictate events in the east and began a long, costly retreat. Weinberg notes that the eastern theatre absorbed the overwhelming share of German military losses, making Soviet endurance and counteroffensive central to Allied victory. The battle became an enduring symbol of resistance and sacrifice, illustrating how a single campaign could decisively alter the strategic balance of a global conflict and reshape its ultimate outcome. \cite{beevor2012} \cite{weinberg1994} \cite{hastings2011}

\subsection{Pearl Harbor and the Globalization of the War}
While the Eastern Front consumed German strength, events in the Pacific transformed the character of the conflict. Japan, pursuing imperial expansion across East Asia, faced mounting friction with the United States over its war in China and its move into Indochina. American economic pressure, including an oil embargo, confronted Japanese leaders with a stark choice between retreat and escalation. Calculating that only a swift, paralyzing blow could secure their ambitions before American power fully mobilized, they planned a surprise attack on the United States Pacific Fleet at Hawaii.

On 7 December 1941 Japanese carrier aircraft struck Pearl Harbor, sinking or damaging numerous battleships and killing thousands of servicemen. The attack achieved tactical surprise but failed to destroy the American aircraft carriers, which were absent, or to cripple repair and fuel facilities. Strategically, it proved a grave miscalculation, uniting a previously divided American public behind total war. President Roosevelt's call to arms transformed isolationist sentiment into determined national resolve, committing the immense industrial and military resources of the United States to the defeat of Japan.

Days later Hitler declared war on the United States, fulfilling commitments to Japan and reflecting his contempt for American military capacity. This decision, though not strictly required by treaty, fused the European and Pacific wars into a single global struggle. The United States now joined Britain and the Soviet Union in a grand coalition spanning every major theatre. Weinberg stresses that the entry of American power, combined with Soviet endurance, established a material and strategic foundation that the Axis could not ultimately match across the breadth of the conflict.

The early months following Pearl Harbor brought a string of Japanese victories across the Pacific and Southeast Asia, including the fall of Singapore and the Philippines. Yet these conquests stretched Japanese forces across vast distances while American industry mobilized for sustained war. The carrier battles at Coral Sea and Midway in 1942 halted Japanese expansion and inflicted irreplaceable losses on the Imperial Navy. From that point the United States pursued a relentless offensive, leveraging its growing fleet and productive capacity to roll back Japanese gains across the ocean.

Hastings emphasizes that American entry altered not only the balance of forces but the entire logic of Allied strategy. The United States could supply allies, sustain multiple offensives, and absorb losses on a scale impossible for the Axis. Coordination among the Allied powers, despite tensions over priorities and timing, channelled this immense potential toward coherent objectives. Pearl Harbor thus marked the moment the war became truly global, binding distant theatres into an interconnected struggle whose outcome would increasingly reflect the disparity in resources between the contending coalitions. \cite{hastings2011} \cite{weinberg1994} \cite{beevor2012}

\bigskip
\begin{otherlanguage}{hebrew}
\subsection*{סיכום הפרק}

כיתור וכניעת הארמייה השישית הגרמנית בסטלינגרד בתחילת {\beginL 1943\endL} העבירו את היוזמה האסטרטגית לידי ברית המועצות. התקפת יפן על פרל הארבור בדצמבר {\beginL 1941\endL} הביאה את העוצמה האמריקנית למלחמה והפכה אותה לעימות עולמי.

\end{otherlanguage}

\clearpage
